\documentclass[11pt]{article}
\usepackage[margin=1in]{geometry}
\usepackage[T1]{fontenc}
\usepackage[utf8]{inputenc}
\usepackage{booktabs}
\usepackage{hyperref}

\title{GPQA Calibration Under the Repaired JSON MC Contract}
\date{2026-03-16 UTC}

\begin{document}
\maketitle

\section*{Bottom Line}
The GPQA evaluation path is back in the expected range once the run uses the repaired
JSON-answer contract together with a benchmark-style decode policy instead of the
bounded 4096-token pilot. The full \texttt{GPQA Diamond} pass landed at
\textbf{78 / 198 = 39.39\%}, which is close to Qwen's own reported
\textbf{40.1} for \texttt{Qwen3-1.7B} in thinking / BF16 mode.\footnote{\url{https://huggingface.co/Qwen/Qwen3-GPTQ-Int8}}

\section*{Run Configuration}
I reran the full local \texttt{gpqa\_diamond.csv} mirror under the repaired rollout
collector contract with:
\begin{itemize}
\item model \texttt{Qwen/Qwen3-1.7B}
\item repaired terminal JSON answer contract
\item \texttt{temperature=0.6}
\item \texttt{num\_generations=1}
\item \texttt{max\_tokens=38912}
\item \texttt{max\_model\_len=40960}
\item single-GPU BF16 decode on \texttt{wth-gpu-01}
\end{itemize}

These settings intentionally target the benchmark-style configuration Qwen recommends
for thinking mode: \texttt{temperature=0.6}, \texttt{top\_p=0.95},
\texttt{top\_k=20}, and a long output budget for difficult problems.\footnote{\url{https://huggingface.co/Qwen/Qwen3-1.7B}}

\section*{Result}
\begin{center}
\begin{tabular}{lrr}
\toprule
Quantity & Value & Note \\
\midrule
Correct & 78 / 198 & 39.39\% accuracy \\
Looped & 4 / 198 & 2.02\% \\
Max-length hits & 0 / 198 & 0.00\% \\
Average generation length & 7302.73 & tokens \\
Average looped generation length & 12365.75 & tokens \\
Average first-loop prefix length & 7199.25 & tokens \\
\bottomrule
\end{tabular}
\end{center}

\section*{Why This Matters}
The earlier bounded parser pilot was never benchmark-comparable for GPQA because it
used \texttt{max\_tokens=4096}; under that pilot contract the corrected JSON parser
still reached only \texttt{33 / 160 = 20.63\%} on the 16-prompt / 160-rollout slice,
with \texttt{101 / 160} responses ending by \texttt{length}. By contrast, the full
calibration run above eliminates max-length failures entirely and lands within about
\texttt{0.7} percentage points of Qwen's own published GPQA number.

\section*{Interpretation}
\begin{itemize}
\item The repaired JSON-answer contract is good enough to keep.
\item The 4096-token truncation concern belongs to the bounded parser pilots, not to
the benchmark-style GPQA evaluation path.
\item For GPQA, the evaluation now looks benchmark-comparable rather than chance-level.
\item The remaining distinction to preserve is methodological: this benchmark-style
GPQA check used \texttt{num\_generations=1}, while the rollout-stat bundle used
multi-sample settings for loop telemetry.
\end{itemize}

\end{document}
